% =============================================================================
% PREFACIO — Acerca de esta plantilla
% =============================================================================
\chapter*{Acerca de esta plantilla}
\addcontentsline{toc}{chapter}{Acerca de esta plantilla}

\begin{flushright}
\itshape
``El laboratorio computacional es el espacio donde la teoría se encuentra con la simulación.''
\end{flushright}
\bigskip

Esta plantilla \LaTeX\ está diseñada para la elaboración de compendios de prácticas de laboratorio en ciencias, con especial énfasis en física computacional y machine learning. Proporciona una estructura modular, entornos personalizados (definiciones, teoremas, ejemplos, ejercicios computacionales), resaltado de código con la paleta Catppuccin, y un formato orientado al estilo IMRaD (Introducción, Métodos, Resultados, Discusión).

\subsection*{Estructura del compendio}
El documento principal organiza las prácticas como capítulos independientes. Cada práctica puede contener:
\begin{itemize}
  \item \texttt{abstract} – resumen de la práctica.
  \item \texttt{objetivos} – metas específicas.
  \item \texttt{fundamentos teóricos} – conceptos necesarios.
  \item \texttt{métodos propuestos} – algoritmos descritos sin lenguaje de programación.
  \item \texttt{experimentos numéricos} – lista concreta de tareas a realizar.
  \item \texttt{resultados} – gráficas y tablas.
  \item \texttt{discusión} – análisis e interpretación.
  \item \texttt{conclusiones} – síntesis final.
  \item \texttt{notas bibliográficas} – referencias.
\end{itemize}

\subsection*{Personalización}
Para adaptar esta plantilla a su propio compendio, modifique:
\begin{itemize}
  \item \texttt{metadata.tex} – título, autor, afiliación, palabras clave, etc.
  \item \texttt{frontmatter/titlepage.tex} – portada.
  \item \texttt{config/geometry.tex} – márgenes y diseño de página.
  \item \texttt{config/titles-config.tex} – formato de capítulos y secciones.
\end{itemize}

El código fuente de los ejercicios puede incluirse mediante el entorno \texttt{lstlisting} con los estilos predefinidos para Python, Julia, C++, etc. No se requiere un lenguaje específico; el estudiante puede implementar los algoritmos en el entorno que prefiera.

\vspace{1.2cm}
\begin{flushright}
\textit{Jesús Antonio Chala Casanova}\\[2pt]
\the\year
\end{flushright}